\chapter{Introduction: What We Are Building and Why Gold}
\label{ch:intro}

\section{The goal}

You are going to write, from scratch, a program that solves the quantum
mechanics of a real crystalline solid. The input will be almost embarrassingly
small: the chemical species (gold), the crystal structure (face-centered
cubic), and a lattice constant. From this the program will predict, with no
empirical fitting whatsoever, the electronic band structure, the total energy,
the equilibrium lattice constant, the bulk modulus, the density of states,
a topological invariant of bulk gold, the ingredients of its electrical
conductivity, and its reflectivity spectrum and colour --- quantities you
can check against experiment and against
mature codes such as Quantum ESPRESSO, VASP, or ABINIT.

The framework is \emph{Kohn--Sham density functional theory} (DFT), the
workhorse of modern computational materials science. This document contains
\emph{all} the theory and \emph{all} the algorithms you need, assuming only
undergraduate quantum mechanics: the time-independent Schr\"odinger equation,
Dirac notation, Hermitian operators, the variational principle, the
hydrogen atom, spin-$\tfrac12$ algebra with Pauli matrices, and first-order
perturbation theory --- plus enough Python and NumPy to write and debug a
few hundred lines of array code (the handful of JAX features used are
explained in Chapter~\ref{ch:gpu}). Everything
else --- functional calculus, the Hohenberg--Kohn theorems, exchange and
correlation, Bloch's theorem, pseudopotentials, relativistic corrections,
iterative eigensolvers, self-consistency acceleration --- is developed here.

No program listings appear in this document, by design: writing the code is
your job, and it is how you will learn the method. Every algorithm is instead
given as explicit mathematics plus pseudocode (the handful of JAX functions
you need are named and explained in Chapter~\ref{ch:gpu}), and each
computational chapter ends
with notes on how the algorithm maps onto a massively parallel accelerator
(GPU) under JAX: which loops become batched tensor contractions, which
operations are FFTs, where \texttt{vmap}-style batching applies,
and which numerical choices (block eigensolvers, Fourier-space mixing) exist
\emph{specifically because} they parallelize well.

\section{Why gold is the right ``hard'' first material}

A first DFT code is usually tested on silicon: a light-element semiconductor
where spin can be ignored and a coarse Brillouin-zone sampling suffices. This
guide deliberately aims at something harder, and gold is the ideal choice,
because it is
\emph{physically} demanding while remaining \emph{structurally} trivial:

\begin{enumerate}[leftmargin=2em]
\item \textbf{One atom per primitive cell.} FCC gold has a single atom in the
primitive cell. Your arrays stay small, debugging stays sane, and yet nothing
essential is simplified away.

\item \textbf{It is a metal.} There is no band gap. Occupation numbers are not
simply ``empty'' or ``full''; you must locate a Fermi level, introduce fractional
(smeared) occupations, sample the Brillouin zone densely, and stabilize the
self-consistency loop against ``charge sloshing.'' These are exactly the
algorithmic complications that separate a toy code from a real one
(Chapters~\ref{ch:solids} and \ref{ch:algorithms}).

\item \textbf{The $5d$ electrons make exchange--correlation matter.} Gold's
valence configuration is $5d^{10}6s^1$. The filled, spatially compact $d$
shell sits only $\sim$2--8\,eV below the Fermi level and hybridizes with the
$6s$ band. The position of these $d$ bands --- and hence the computed optical
absorption edge, lattice constant, and cohesive energy --- is sensitive to the
exchange--correlation functional. You will implement both the local density
approximation (LDA) and the PBE generalized gradient approximation and see
their differences directly (Chapter~\ref{ch:xc}).

\item \textbf{Spin--orbit coupling is enormous and observable.} With nuclear
charge $Z=79$, the $5d$ electrons of gold move at an appreciable fraction of
the speed of light near the nucleus. Spin--orbit coupling (SOC) splits the
$d$ bands at the $\Gamma$ point by more than $1$\,eV --- not a small correction
but a dominant feature of the band structure. Famously, relativity is why
gold is yellow while silver (same column, $Z=47$) is white: relativistic
stabilization of the $6s$ band and destabilization of the $5d$ band shift the
interband absorption edge from the ultraviolet (silver, $\approx 3.8$\,eV)
down into the visible ($\approx 2.4$\,eV, blue-green light), so gold absorbs
blue and reflects red and yellow. Your code
will treat SOC non-perturbatively with two-component spinor wavefunctions
(Chapter~\ref{ch:soc}).

\item \textbf{It hides a topological invariant.} Although gold is a metal,
its bands carry a nontrivial $\mathbb{Z}_2$ topological index, and the famous
surface state on Au(111) is its fingerprint. This is not an afterthought
bolted onto the project: the index is a property of the \emph{wavefunctions}
your code produces, and every tool needed to extract it is introduced where
it naturally belongs --- the parity of a zone-boundary state in the very
first band-gap model (Section~\ref{sec:nfe}), inversion and time reversal as
operations on plane-wave coefficients (Section~\ref{sec:symmetry}), and the
fixed-potential band solver (Section~\ref{sec:nscf}).
Chapter~\ref{ch:topology} then only has to assemble them.

\item \textbf{Excellent reference data.} The experimental lattice constant
($a = 4.0782$\,\AA), bulk modulus ($\approx 173$\,GPa; both are quoted at
room temperature, while DFT describes a static lattice at $0$\,K: removing the
zero-point and thermal expansion shrinks $a$ by $\approx 0.015$\,\AA{}, to
$\approx 4.063$\,\AA, and stiffens the lattice: $B \approx 180$\,GPa from the
$0$\,K elastic constants, up to $\approx 196$\,GPa in the static-lattice
estimates used by density-functional benchmarks), the photoemission
$d$-band positions, and published DFT band structures at every level of
theory give you unambiguous targets for verification
(Chapter~\ref{ch:verify}).
\end{enumerate}

\begin{keybox}{The physical system}
FCC gold. Primitive lattice vectors (with $a = 4.0782$\,\AA{}
$= 7.706681$ Bohr; convert with the full $a_0 = 0.529177211$\,\AA{} ---
rounding $a$ to four decimals already shifts the Ewald energy by
$10^{-4}$\,Ha and fails the tests of Chapter~\ref{ch:verify}):
\[
\mathbf{a}_1 = \tfrac{a}{2}(0,1,1),\qquad
\mathbf{a}_2 = \tfrac{a}{2}(1,0,1),\qquad
\mathbf{a}_3 = \tfrac{a}{2}(1,1,0),
\]
one Au atom at the origin. With a scalar-relativistic + SOC pseudopotential
treating $5d^{10}6s^1$ (11 electrons) or $5s^25p^65d^{10}6s^1$ (19 electrons,
``semicore,'' more accurate) explicitly.
\end{keybox}

\section{Units}

We use \textbf{Hartree atomic units} throughout unless stated otherwise:
$\hbar = m_e = e = 4\pi\varepsilon_0 = 1$. Then lengths are in Bohr radii
($a_0 = 0.529177$\,\AA), energies in Hartree ($1\,\mathrm{Ha} = 27.2114$\,eV
$= 2$\,Ry), and the speed of light is $c = 1/\alpha \approx 137.036$, where
$\alpha$ is the fine-structure constant. The one-electron kinetic operator is
$-\tfrac12\nabla^2$ and the Coulomb interaction between electrons is
$1/|\rr_1-\rr_2|$. Atomic units eliminate an entire class of bugs; convert to
eV and \AA{} only at the input/output boundary of the program.

\section{Bird's-eye view of the whole machine}

It helps to see the destination before the journey. The final program is a
fixed-point iteration --- the \emph{self-consistent field} (SCF) loop ---
wrapped around a large eigenvalue problem:

\begin{enumerate}[leftmargin=2em]
\item \textbf{Guess} an electron density $n(\rr)$ (e.g.\ a normalized Gaussian
blob of valence charge on each atom).
\item \textbf{Build the effective potential} $\veff[n](\rr)$: the sum of the
ionic (pseudo)potential, the classical electrostatic (Hartree) potential of
the density itself, and the exchange--correlation potential, which encodes
all quantum many-body effects (Chapters~\ref{ch:hkks}--\ref{ch:xc}).
\item \textbf{Solve} the single-particle Schr\"odinger-like (Kohn--Sham)
equation in this potential at many points $\kk$ of the Brillouin zone, using
a plane-wave basis, obtaining bands $\varepsilon_{n\kk}$ and two-component
spinor wavefunctions $\psi_{n\kk}$ (Chapters~\ref{ch:solids},
\ref{ch:soc}, \ref{ch:hamiltonian}). This is done with an iterative,
GPU-friendly block eigensolver (Chapter~\ref{ch:algorithms}).
\item \textbf{Occupy} the bands according to a smeared Fermi--Dirac-like
distribution, locating the Fermi level so the electron count is correct.
\item \textbf{Recompute the density} from the occupied wavefunctions.
\item \textbf{Mix} the new density with the old ones (Pulay/Anderson mixing
with Kerker preconditioning --- what keeps metallic cells stable) and go to step 2 until
the density stops changing.
\item \textbf{Evaluate} the total energy, band structure, density of
states, and wavefunction-based quantities (parities, Berry phases,
velocity matrix elements for transport and optics). (Forces
on the atoms vanish by symmetry with one atom per cell, so they are not
needed here; the equilibrium geometry is found from the energy as a function
of the lattice constant.)
\end{enumerate}

Almost all floating-point work lives in step 3, and there it lives inside two
kernels: \emph{batched 3D FFTs} and \emph{batched dense linear algebra}. Both
are the native food of GPUs. That is not a coincidence --- the plane-wave
formulation of DFT was co-evolved with vector and parallel hardware, and the
algorithm choices in this document (LOBPCG over sequential Davidson,
Fourier-space Kerker mixing, Monkhorst--Pack batching over $\kk$-points) are
made with your JAX target explicitly in mind.

\section{How to use this document}

Read Chapters~\ref{ch:manybody}--\ref{ch:soc} for the physics, with pen and
paper; they contain no programming. Chapters~\ref{ch:hamiltonian} and
\ref{ch:algorithms} are the specification of the program, and
Chapter~\ref{ch:gpu} translates it into JAX. Then \emph{build in the order
of Chapter~\ref{ch:verify}}: every stage there has a number that your code
must reproduce before you move on, so that at no point are you debugging
more than one new ingredient. Three conventions will save you days:
\begin{itemize}[leftmargin=2em]
\item Boxes with a blue-grey frame summarize what must be remembered; boxes
with a green frame contain implementation advice for JAX.
\item Every formula that you will type into the computer comes with a
\emph{numerical anchor} --- a value to compare against --- either next to it
or in Chapter~\ref{ch:verify}. Use them all; a wrong constant is invisible
in a band plot.
\item The folder \texttt{verification/} distributed with this document
holds the scripts that produced every reference number, including an
independent JAX implementation written from this text alone. Treat it as
the answer key: consult it only after your own attempt.
\end{itemize}

\section{Chapter map}

\begin{itemize}[leftmargin=2em]
\item \textbf{Ch.~\ref{ch:manybody}} --- the exact many-electron problem and
why it is intractable; the Born--Oppenheimer approximation; the density as
the central variable.
\item \textbf{Ch.~\ref{ch:hkks}} --- functional calculus from zero; the
Hohenberg--Kohn theorems with proofs; the Kohn--Sham construction and a full
derivation of the Kohn--Sham equations.
\item \textbf{Ch.~\ref{ch:xc}} --- exchange and correlation: the homogeneous
electron gas, derivation of LDA exchange, the PW92 correlation
parametrization, the PBE GGA, and the (non)collinear spin generalization
needed once SOC is on.
\item \textbf{Ch.~\ref{ch:solids}} --- crystals: Bravais and reciprocal
lattices, Bloch's theorem (proved), plane-wave basis sets and cutoffs, the
two-plane-wave model of a band gap, Brillouin-zone integration,
Monkhorst--Pack meshes, metallic occupations and smearing, the Fermi level.
\item \textbf{Ch.~\ref{ch:pseudo}} --- pseudopotentials: why all-electron
plane-wave calculations are hopeless, norm conservation, the
Kleinman--Bylander separable form, and the analytic relativistic HGH
pseudopotentials you will use for gold.
\item \textbf{Ch.~\ref{ch:soc}} --- relativity: from the Dirac equation to
the Pauli Hamiltonian, the physical origin of SOC, $j$-dependent
pseudopotentials, two-component spinor Kohn--Sham theory, time-reversal and
inversion symmetry as explicit operations on your coefficient arrays, and
what SOC does to gold's bands.
\item \textbf{Ch.~\ref{ch:hamiltonian}} --- the full Hamiltonian in the
plane-wave basis, term by term: kinetic, local via FFT, nonlocal via
projectors, Hartree in reciprocal space, Ewald summation for the ion--ion
energy, and the total-energy expression.
\item \textbf{Ch.~\ref{ch:algorithms}} --- the algorithms: basis
construction, matrix-free Hamiltonian application, preconditioned block
eigensolvers (LOBPCG and blocked Davidson), subspace diagonalization,
Fermi-level bisection, density construction, Kerker-preconditioned Pulay
mixing, the complete SCF driver, and the fixed-potential solver used for
all post-processing --- all as pseudocode.
\item \textbf{Ch.~\ref{ch:gpu}} --- mapping everything onto JAX/GPU: tensor
layouts, batching over $\kk$ and bands, FFT strategy, mixed precision,
JIT-compatible control flow, and memory budgeting.
\item \textbf{Ch.~\ref{ch:topology}} --- the topology of gold's bands:
Berry phases and Wannier charge centers from scratch, the $\mathbb{Z}_2$
invariant as a statement about how those centers flow, its evaluation from
wavefunction parities and by Wilson loops, with pseudocode.
\item \textbf{Ch.~\ref{ch:conductivity}} --- electrical conductivity: the
velocity operator, the Drude plasma frequency from the Fermi surface, and
an honest account of what a ground-state calculation cannot give (the
scattering time).
\item \textbf{Ch.~\ref{ch:colour}} --- the colour of gold: the dielectric
function from interband transitions, the reflectivity, and the conversion
of a reflectance spectrum into the colour a human eye sees.
\item \textbf{Ch.~\ref{ch:verify}} --- a staged verification plan with
quantitative targets, from the free-electron gas to the SOC-split $d$ bands
of gold and its topological index, including reference numbers from an
independent implementation of exactly the equations in this document.
\item \textbf{Appendix~\ref{app:q19}} --- the semicore (19-electron)
pseudopotential parameters.
\end{itemize}

\section{Notation at a glance}
\label{sec:notation}

The same few indices appear in every chapter; fix them now.
\begin{center}
\begin{tabular}{@{}lp{0.66\textwidth}@{}}
\toprule
Symbol & Meaning\\
\midrule
$\rr$, $\RR$ & position in space; lattice translation (or atomic position
$\RR_a$)\\
$\kk$, $\GG$ & wavevector in the Brillouin zone; reciprocal lattice vector\\
$n$ (as an index), $\kk$ & band index and wavevector of a Bloch state
$\psi_{n\kk}$\\
$n(\rr)$, $\mathbf m(\rr)$ & electron density; magnetization density\\
$s$, $\alpha$, $\beta$ $\in\{\uparrow,\downarrow\}$ & spinor component\\
$c^{s}_{n\kk}(\GG)$ & plane-wave coefficients of a spinor Bloch state
(the program's central array)\\
$f_{n\kk} \in [0,1]$ & occupation of one spinor state (one electron at
most)\\
$w_{\kk}$ & Brillouin-zone weight of a $\kk$-point, $\sum_{\kk} w_{\kk} = 1$\\
$\varepsilon_{n\kk}$, $\mu$ & Kohn--Sham eigenvalue; Fermi level (chemical
potential)\\
$\Omega$, $N_e$ & volume of the primitive cell; number of valence electrons
per cell\\
$\ell$, $m$, $i$ & angular momentum, its projection, and radial index of a
pseudopotential projector\\
$E_{\mathrm{cut}}$, $\sigma$ & plane-wave cutoff; smearing width\\
\bottomrule
\end{tabular}
\end{center}
A hat ($\hat H$) marks an operator; a tilde ($\tilde p$, $\tilde V$) marks
a Fourier- or Bessel-transformed function. ``Pair'' always means a
Kramers pair of spinor bands (two degenerate bands; Section~\ref{sec:symmetry}).
A few letters are unavoidably reused, each time with the meaning standard
in that field, and the context always says which is meant: $\sigma$ is the
smearing width, but $\bm\sigma$ the Pauli matrices and, in
Chapter~\ref{ch:conductivity}, the conductivity; $s$ is a spinor index, but
also the smearing entropy $s(x)$, the reduced density gradient of PBE and a
Gaussian width; $\mu$ is the Fermi level, but also a PBE constant and the
index of a real spherical harmonic; $\alpha,\beta$ are spinor indices, but
also the fine-structure constant, the PW92 and PBE constants and, in
Chapter~\ref{ch:conductivity}, Cartesian directions; and $A$, $B$, $K$, $P$
each have several meanings (lattice matrices, PW92/PBE constants, mixing
amplitude; preconditioner, complex conjugation, Kerker operator, the
$K$ point; parity operator, projections, an eigensolver block).
